Sincere thanks are extended to Vincenzo Messina for his supervision throughout this thesis. The discussions were particularly valuable when the work moved from running architecture cases to deciding which comparisons were scientifically meaningful. He provided the original communication-model basis, supported the integration with TAT-C, and helped define the questions used to evaluate the final campaigns.

The Chair of Spacecraft Systems at the Technical University of Munich is also thanked for providing the academic and computing environment used for this work. Gratitude is likewise extended to family and friends for their patience and support during the long simulation and writing phases.

OpenAI's ChatGPT, including the Codex environment, was used as an assistive tool during this thesis. Its use included support with Python code development and debugging, preparation of simulation and post-processing workflows, generation and refinement of plots and tables, language correction, and restructuring of draft sections in response to supervisory feedback. The research questions, modelling assumptions, input data, simulation design, selection of comparisons, validation decisions, and scientific interpretation remained the responsibility of the author. All AI-assisted code, calculations, figures, and text were reviewed against the simulation outputs, source literature, and supervisory feedback before inclusion. ChatGPT was not treated as a scientific source and is not used to substantiate technical claims.
